Hij werd langzaam vager.
Dromend werd de jager
een norse hoorndrager

met een puntige staart.
Hij kreeg een stoppelbaard.
Een drietand was zijn zwaard.

Hij werd enorm gespierd.
Met een six-pack versierd.
Hij leek ongemanierd

met zijn zwarte vleugels.
In zijn handen beugels,
gekoppeld aan teugels,

waarmee hij twee vrouwen
aan zich bond met touwen.
Ontbloot in zijn klauwen.

Zijn ogen vol terreur
waren felwit van kleur.
Hij verspreidde een geur,

als die van Raspoetin.
Emmy kreeg meteen zin.
De feromonen in

haar neus lokten genoeg.
De macht die hij uitdroeg,
tintelde haar schaamvoeg.

Ze werd zeer genegen
om zich aan het wezen
ongeremd te geven.

Mijmerend en besmeurd
werd Emmy meegesleurd.
Ze hintte op haar beurt.

Ze wilde haar portie.
Kon de jager Emmy
helpen haar fantasie

heerlijk uit te leven?
Het monster beleven?
De jager kon geven:

“Je bent een geil sletje,
in voor een verzetje
in elk vacant bedje.

Een hoer in vaktermen.
Dat gaat je beschermen.
Ik profeteer zwermen,

die je gaan opzoeken.
Talloze aanzoeken,
die je zult vervloeken,

omdat je mijn hoer bent.
En als je mij goed kent
dan ben ik niet de vent,

die voor je gaat zorgen.
Maar je zult vast morgen
je lust gaan waarborgen.

Wat jij wilt boeit me niet.
Ik wil een slaafse griet,
waarop ik vaak geniet.

Blinde overgave!
Vrije offergave!”
Ze liet hem doordraven,

maar werd wel aandachtig
toen hij plots schaapachtig
haar kende waarachtig.

“Jij lijkt net op Emma
uit de hoerencentra
waar ik dikwijls heen ga.

Je voornaam lijkt er op.
Net als je rode kop.
Ik vond haar lichaam top.

Een reisvriend was eens mee,
deed in haar zijn entree
en neukte haar tevree.

Ze was vermomd dit keer.
Een woedende landheer
stak haar ineens dood neer,

onder bar woordgebruik,
toen hij haar zwarte pruik,
wit bevlekt door misbruik,

van Emma’s hoofdje trok.
De landheer was in schok
door de vrucht van zijn gok.

Ze bleek ook gemaskeerd,
terwijl hij gedrogeerd
haar had gepenetreerd

de zwoele nacht ervoor.
De landheer had niet door
dat hij zijn eer verloor.

Ze deed het werk als snol
kennelijk voor de lol.”
Emmy’s gemoed schoot vol,

al zag hij niet haar pijn.
Ze herschreef toen haarfijn
haar zusters verhaallijn.

Emmy achterhaalde
haar list en vertaalde
wat Emma verhaalde

als onbegrensde lust.
Emmy was toegerust
met wat Emma belust.

Dus met haar in de buurt
is er niemand die gluurt
naar wat Emma toestuurt.

Emma was vaak gegriefd,
want ze was vaak verliefd.
“Emmy, pak alsjeblieft

niet mannen van me af,
zodat ik géén aanschaf.
Pak ze niet als maatstaf

voor jouw uitstekendheid,
superioriteit,
wat betreft je schoonheid.”,

waren ongeschreven
woorden die aan leven
met haar zusje kleven.

Emmy was ook te vroom
om Emma zonder schroom
haar erotische droom

te laten uitleven.
Al duurde het even,
ze zou haar aangeven,
al nam het diens leven.

In haar schuldbewustheid
wist ze met zekerheid:
"Ik ben mijn zuster kwijt”.

Een affect verwijst doorgaans naar een intense of kortdurende emotionele reactie die vaak niet bewust wordt gecontroleerd. Het hoeft niet per se iets te maken te hebben met het verleden of met vorige levens; het is eerder een directe, spontane uiting van emotionele energie, zoals blijdschap, woede of verdriet.

De mythe van Inanna's afdaling naar de onderwereld is een van de oudste literaire werken die bekend zijn uit het oude Soemerië. Het verhaal beschrijft de beproevingen en transformaties van de godin Inanna, ook bekend als Ishtar, tijdens haar reis naar de onderwereld. Deze mythe kan worden gezien als een archetypische reis van de ziel, waarin thema's zoals sterven en wedergeboorte, opoffering, en de confrontatie met het duistere zelf worden verkend.
Samenvatting van de mythe:
Inanna, de godin van liefde, vruchtbaarheid en oorlog, besluit om af te dalen naar de onderwereld, het domein van haar zuster Ereshkigal, de koningin van de dood. Voordat ze vertrekt, instrueert ze haar dienares Ninshubur om hulp te zoeken bij de goden als ze niet terugkeert. Inanna draagt al haar machtige attributen en symbolen, zoals haar kroon, sieraden, en koninklijke gewaden.
Bij haar aankomst aan de poorten van de onderwereld, wordt Inanna door de poortwachter Neti gestopt. Hij informeert Ereshkigal over haar komst, en Ereshkigal besluit dat Inanna de onderwereld mag binnengaan, maar alleen als ze zich van haar koninklijke attributen ontdoet. Inanna moet door zeven poorten gaan, en bij elke poort moet ze één van haar machtssymbolen afgeven, totdat ze uiteindelijk naakt en kwetsbaar is. Dit symboliseert het loslaten van haar ego, macht en identiteit.
Wanneer ze eindelijk voor Ereshkigal staat, is ze volledig ontdaan van haar wereldlijke kracht en wordt ze ter dood gebracht. Haar lichaam wordt opgehangen aan een haak, waar het drie dagen en drie nachten blijft hangen. Deze fase vertegenwoordigt een symbolische dood en een staat van volledige overgave en ontbering.
Ninshubur, trouw aan haar belofte, vraagt de goden om hulp. Uiteindelijk wordt Enki, de god van wijsheid, bewogen om in te grijpen. Hij creëert twee wezens van zijn nagels en stuurt hen naar de onderwereld met het voedsel en water des levens. Ze sympathiseren met het lijden van Ereshkigal en verkrijgen zo de toestemming om Inanna’s lichaam te besprenkelen, waardoor ze weer tot leven komt.
Hoewel Inanna uit de onderwereld mag terugkeren, kan niemand de onderwereld verlaten zonder een vervanger achter te laten. Ze keert terug naar de bovenaardse wereld, maar uiteindelijk moet haar geliefde echtgenoot, Dumuzi, de straf ondergaan en in haar plaats afdalen naar de onderwereld.
Symbolische Interpretatie:
Afdaling en Ontbering: Inanna's afdaling kan gezien worden als een archetypische 'weg van de beproevingen' waarin de held(in) zichzelf verliest om een dieper, spiritueel inzicht te verkrijgen. De zeven poorten symboliseren de lagen van ego die moeten worden afgeworpen om innerlijke transformatie te bereiken.
Dood en Wedergeboorte: Haar dood en herleving vertegenwoordigen de cyclische aard van spirituele transformatie, waarbij de dood niet het einde is, maar een noodzakelijk proces van loslaten en wedergeboorte.
Confrontatie met het Schaduwzelf: De ontmoeting met Ereshkigal, haar donkere zuster, kan worden gezien als een confrontatie met het schaduwzelf, de onderdrukte of ongeaccepteerde aspecten van de psyche.
Offer en Wederkeer: Het offer van Dumuzi vertegenwoordigt de balans en de onvermijdelijke noodzaak van offer en wederkeer. De les hier is dat er voor elke wedergeboorte een offer nodig is.
De mythe van Inanna's afdaling naar de onderwereld biedt diepe inzichten in de menselijke en spirituele zoektocht, waarin beproevingen dienen als katalysatoren voor groei en transformatie.
